% source: An algebraic approach to the Riemann-Roch theorem and the arithmetic theory of function fields (Athanasios Papaioannou), http://www.fen.bilkent.edu.tr/~franz/mat/Riemann-Roch.pdf
% pdf_page: 0009
% retrieved: 2026-07-26
% transcribed_by: page-transcriber (AI vision, no OCR)

\documentclass[11pt]{article}
\usepackage{amsmath,amssymb,amsthm}

% Small-caps theorem style matching the source.
\newtheoremstyle{scplain}{\topsep}{\topsep}{\itshape}{0pt}{\scshape}{.}{ }{\thmname{#1} \thmnote{#3}}
\theoremstyle{scplain}
\newtheorem*{corollary}{Corollary}
\newtheorem*{theorem}{Theorem}
\newtheorem*{definition}{Definition}

% Small-caps "Proof." to match the source.
\makeatletter
\renewenvironment{proof}[1][\proofname]{%
  \par\pushQED{\qed}\normalfont\topsep6\p@\@plus6\p@\relax
  \trivlist\item[\hskip\labelsep\scshape #1\@addpunct{.}]\ignorespaces}%
  {\popQED\endtrivlist\@endpefalse}
\makeatother
\renewcommand{\qedsymbol}{$\square$}

\DeclareMathOperator{\Div}{Div}

\begin{document}

\begin{corollary}[1.16]
\begin{enumerate}
\item[(a)] The degrees and the dimensions of any two equivalent divisors
coincide.

\item[(b)] If $\deg D < 0$, then $\dim D = 0$.

\item[(c)] For a divisor $D$ of degree $0$, the following assertions are
equivalent:
  \begin{enumerate}
  \item[(i)] $D$ is principal
  \item[(ii)] $\dim D \geq 1$
  \item[(iii)] $\dim D = 1$.
  \end{enumerate}
\end{enumerate}
\end{corollary}

Notice now that in Lemma 1.14 we saw that $\dim D \leq 1 + \deg D$ for any
effective divisor $D$. However, this holds for all divisors $D$ with
$\deg D \geq 0$. In order to verify this, assume that $\dim D > 0$, in which
case $D \sim D'$ for some effective divisor $D'$, the result now follows
immediately, by the proved equalities $\deg D = \deg D'$ and
$\dim D = \dim D'$.

The desired theorem is now attainable. We have

\begin{theorem}[1.17]
Any function field has finite genus. More precisely, there is a constant
$\gamma \in \mathbb{Z}$ that only depends on $\mathbb{F}/\mathbb{K}$ and is such
that
\[
  \deg D - \dim D \leq \gamma,
\]
for any $D \in \Div(\mathbb{F})$.
\end{theorem}

\begin{proof}
First note that $D_1 \leq D_2$ implies that
$\deg D_1 - \dim D_1 \leq \deg D_2 - \dim D_2$, by Lemma 1.13. Now fix a
non-constant element $x \in \mathbb{F}$ and consider the specific divisor
$B = (x)_\infty$. As mentioned before, there exists an effective divisor $C$
such that $\dim(kB+C) \geq (k+1)\deg B$, for all non-negative integers $k$. On
the other hand, $\dim(kB+C) \leq \dim(kB)+\dim C$, by lemma 1.13; the above
inequalities combined yield
\[
  \dim(kB) \geq (k+1)\deg B - \deg C
    = \deg(kB) + ([\mathbb{F} : \mathbb{K}(x)] - \deg C),
\]
hence
\[
  \deg(kB) - \dim(kB) \leq \gamma,
\]
for all $k \geq 0$ and
$\gamma = [\mathbb{F} : \mathbb{K}(x)] - \deg C$. We would like to show that
the above displayed equation continues to hold when we substitute $kB$ by any
$A \in \Div(\mathbb{F})$.

For the above statement to hold, it clearly suffices to prove the following
claim: \emph{given a divisor $A \in \Div(\mathbb{F})$, there exist divisors $A_1$
and $D$ and a positive integer $k$ such that $A \leq A_1$, $A_1 \sim D$ and
$D < kB$} (this statement is enough to finish off the proof, because then
$\deg A - \dim A \leq \deg A_1 - \dim A_1 = \deg D - \dim D
\leq \deg(kB)-\dim(kB) \leq \gamma$). To prove the claim, pick an effective
divisor $A_1$ such that $A_1 \geq A$. Then
$\dim(kB-A_1) \geq \dim(kB)-\deg(A_1) \geq \deg(kB)-\gamma-\deg A_1 > 0$, for
$k \in \mathbb{N}$ sufficiently large. Therefore there is some element
$z \in L(kB-A_1)$ and so we may set $D = A_1-(z)_0$ for a fixed such $z$. Note
that $A_1 \sim D$ and $D \leq A_1-(A_1-kB)=kB$, as desired. This completes the
proof.
\end{proof}

\section*{1.4\quad The Proof of the Riemann-Roch Theorem.}

We will now proceed to prove Roch's generalisation of Riemann's inequality,
which is known as the Riemann-Roch theorem. Before stating it, we will need to
introduce the technical objects that will be used in its proof. The first
important concept is that of an adele.

\begin{definition}[1.11]
An adele of a function field $\mathbb{F}/\mathbb{K}$ is a mapping
$\alpha : \mathbf{P}_{\mathbb{F}} \longrightarrow \mathbb{F}$ defined by
$P \mapsto \alpha_P$ such that $\alpha_P \in \mathcal{O}_P$ for almost all
places $P$ of $\mathbb{F}/\mathbb{K}$. We may regard an adele as an element of
$\prod P$, where the product ranges over all places
$P \in \mathbf{P}_{\mathbb{F}}$.
\end{definition}

The associated space $A_{\mathbb{F}}$ that contains all the adeles of
$\mathbb{F}/\mathbb{K}$ is called the \emph{adele space} and is obviously a
$\mathbb{K}$-vector space, with the algebraic structure inherited from the
identification of $A_{\mathbb{F}}$ with
$\prod_{P \in \mathbf{P}_{\mathbb{F}}} P$; in fact, $A_{\mathbb{F}}$ can be even
regarded as a $\mathbb{K}$-algebra, though this observation will never be used
henceforth.

Given an element $x \in \mathbb{F}$, the \emph{principal adele of $x$} is the
adele whose components are all equal to $x$; this definition agrees with the
conditions imposed on adeles, since any $x \in \mathbb{F}$ has finitely many
zeroes and poles, as we showed earlier. The above construction yields an
embedding $\mathbb{F} \hookrightarrow A_{\mathbb{F}}$, and the valuations $v_P$
of $\mathbb{F}/\mathbb{K}$

\begin{center}9\end{center}

\end{document}
